

\begin{frame}
\frametitle{Hiérarchie des mémoires: en centralisé}


\figSlideWithSize{hiermem}{8cm}

\vfill

Classique: un SGBD place en mémoire les données les plus utilisées
\end{frame}


\begin{frame}
\frametitle{Hiérarchie des mémoires: dans le \textit{cloud}}

Différence essentielle: le réseau (hiérarchique)

\vfill

\figSlideWithSize{hiermem-cloud}{10cm}

\end{frame}

\begin{frame}
\frametitle{Performances\,: ordres de grandeur}

\begin{center}

\begin{small}

\begin{tabular}{l|p{2cm}|p{4cm}|p{3cm}}
{\bf Type} & \textbf{Taille} & {\bf Latence} & \textbf{Débit} \\
\hline
RAM & O(10 GOs) & $\approx 10^{-8} \mathrm{s}$ (10 nanosec.); & $\approx 10 G0/s$ \\ 
SSD & O(100 GOs) & $\approx 10^{-4} \mathrm{s}$ (0,1 millisec.); & Qq GO/s \\ 
Disque & O(1 TOs) & $\approx 10^{-3}\mathrm{s}$ (10 millisec.);&  100 MO/s \\ 
\hline
\end{tabular}

\end{small}
\end{center}

\vfill

Pour le réseau: 
\begin{redbullet}
  \item compter 10 Gbits/s (intra)-baie; 
  \item diviser par 5 ou 10 pour le débit inter-baies.
\end{redbullet}

\vfill

\begin{remblock}{Retenir}
\begin{redbullet}
  \item Accès (aléatoire) au disque\,: très lent par rapport à un accès en RAM.
  \item Débit du disque\,: faible (100 MO/s au mieux)
  \item SSD: 100 fois plus rapide qu'un disque (latence), 10 fois plus en débit
\end{redbullet}
\end{remblock}

\end{frame}


\begin{frame}
\frametitle{Performances dans un \textit{cloud}}

\begin{small}
\begin{tabular}{l|p{1cm}|p{1cm}|p{1cm}|p{1cm}|p{1cm}|p{1cm}}
{\bf Type} & \textbf{RAM locale} & \textbf{Disque local} & \textbf{RAM baie}  & \textbf{Disque baie} 
    & \textbf{RAM cloud}  & \textbf{Disque cloud}   \\
\hline
Latence ($\mu$s) & 0,1 & 10 000  & 300  & 10 000  & 500   &  10 000\\
Débit (MO/s) & 10 000 & 100  & 125  & 100  & 25   & 20\\
\hline
\end{tabular}

\end{small}


\begin{remblock}{Localité des données (\textit{data locality})}
Toujours tenter de traiter les données \alert{localement}, par le CPU 
du serveur stockant les données.
\begin{redbullet}
  \item Pas toujours possible, mais \emph{best effort} pour des systèmes comme Hadoop.
  \item \alert{Implique un déplacement des traitements vers les données} (en analytique)
 \end{redbullet}
\end{remblock}

\vfill

\begin{defblock}{Essentiel}
\begin{redbullet}
  \item Systèmes temps réel: \alert{placer les données en RAM} (latence reste faible).
  \item Systèmes analytiques: \alert{placer les données séquentiellement sur le disque} 
\end{redbullet}
\end{defblock}


\end{frame}

\end{document}


